\documentclass[11pt]{article}
\usepackage[margin=1in]{geometry}
\usepackage{amsmath,amssymb,amsthm}
\usepackage{mathtools}
\usepackage[numbers,square]{natbib}
\usepackage{booktabs}
\usepackage{hyperref}
\hypersetup{hypertexnames=false}
\usepackage{cleveref}

% Theorem environments
\theoremstyle{definition}
\newtheorem{definition}{Definition}[section]
\theoremstyle{plain}
\newtheorem{theorem}{Theorem}[section]
\newtheorem{proposition}[theorem]{Proposition}
\newtheorem{corollary}[theorem]{Corollary}
\newtheorem{lemma}[theorem]{Lemma}
\theoremstyle{remark}
\newtheorem{remark}{Remark}[section]
\newtheorem{example}{Example}[section]

\newcommand{\enc}{\mathrm{enc}}
\newcommand{\dec}{\mathrm{dec}}
\newcommand{\fhat}{\hat{f}}
% Cipher functor: \cipher{X} = C(X) when secret is implicit;
% \cipherS{X}{s} = C_s(X) when secret is explicit.
% Same notation applies to objects (latent types) and morphisms (latent functions).
\newcommand{\cipher}[1]{\mathsf{C}(#1)}
\newcommand{\cipherS}[2]{\mathsf{C}_{#2}(#1)}
\newcommand{\orbitF}{\mathrm{orbit}_F}
\newcommand{\B}{\{0,1\}}

\title{Algebraic Cipher Types: Confidentiality Trade-offs\\in Type Constructors over Trapdoor Computing}
\author{Alexander Towell\\\texttt{lex@metafunctor.com}}
\date{March 2026}

\begin{document}
\maketitle

\begin{abstract}
Cipher types form an algebra: starting from primitive ciphered values,
the type constructors of ordinary programming (products, sums,
exponentials) build complex cipher types from simpler ones, with each
constructor inducing a specific confidentiality cost absent in
plaintext programming.
We characterize the algebra constructor-by-constructor: products
leak correlations, sums force a categorical impossibility (tag
hiding and untrusted pattern matching cannot coexist), and the
cipher exponential is the cipher map abstraction itself, closing
the algebra back on the trapdoor framework.
A uniform information-theoretic bound holds across the algebra:
the conditional entropy of the latent value given the untrusted
machine's view is at least $H(X) - \log_2 |\mathrm{orbit}|$, where
the orbit is the set of cipher values reachable under the operations
the untrusted machine holds.
A typed-composition discipline turns this bound into a design-time
budget: the type of a cipher program determines how deeply its
orbit can grow.
We realize cipher programs as Python expression-tree decompositions
with cipher-map cut points (a \texttt{@cipher\_node} decorator
marks the cipher subtrees; unmarked nodes run as plain code over
opaque bit strings) and evaluate the cipher Boolean specialization
on a 20~Newsgroups Boolean search task, where the cut-point
discipline measurably bounds confidentiality loss.
\end{abstract}


%% ====================================================================
\section{Introduction}
\label{sec:intro}
%% ====================================================================

Trapdoor computing is a paradigm for evaluating functions on data whose
meaning is hidden behind a one-way trapdoor (a cryptographic hash).
A trusted machine encodes plaintext values into cipher values---opaque bit
strings---and hands them, together with cipher maps (total functions on
bit strings), to an untrusted machine.
The untrusted machine evaluates the cipher maps and returns results.
Without the trapdoor (the secret seed), it cannot decode the results or
determine what function it computed.
The formal framework of cipher maps, their four defining properties (totality,
representation uniformity, correctness, composability), and three concrete
constructions (HashSet, entropy map, trapdoor boolean algebra) are developed
in~\cite{towell2026cipher}.

Ordinary programming languages build complex types from simpler ones
using a small set of type constructors: products ($A \times B$, pairs and
records), sums ($A + B$, tagged unions and variants), and exponentials
($A \to B$, functions).
In plaintext computing, these constructors are transparent: a product
exposes both components, a sum carries a tag that enables pattern matching,
and a function is called by providing an argument.
When values are hidden behind a trapdoor, each of these operations
introduces a confidentiality cost that has no plaintext analogue.

This paper identifies and formalizes those costs.
The main contributions are:

\begin{enumerate}
\item \textbf{An algebra of cipher types} (\Cref{sec:types}).
  We survey the standard type constructors over cipher values
  (void, unit, product, sum, exponential) and identify the
  confidentiality cost each one introduces.
  Products leak correlations under component-wise encoding
  (\Cref{prop:product-tradeoff}); sums force an impossibility
  between tag hiding and untrusted pattern matching
  (\Cref{thm:sum-impossibility}); the cipher exponential is
  exactly the cipher map abstraction
  (\Cref{prop:exponential-identity}), closing the algebra back on
  the trapdoor framework, with its own granularity trade-off
  (\Cref{prop:exponential-tradeoff}).
  The encoding-granularity knob (joint vs.\ component-wise) runs
  uniformly through the algebra; the costs differ by constructor.

\item \textbf{A uniform information-theoretic bound} (\Cref{sec:orbit}).
  We define the orbit closure of a cipher value under the
  operations available to the untrusted machine
  (\Cref{def:orbit}), prove monotonicity
  (\Cref{thm:monotonicity}), and derive an entropy-form
  confidentiality bound
  $H(X \mid \mathcal{V}_F(c)) \geq H(X) - \log_2 |\mathrm{orbit}_F(c)|$
  (\Cref{thm:confidentiality-bound}) that is independent of the
  underlying cipher-map construction.
  A typed-composition discipline (\Cref{sec:typed-chains}) converts
  the bound into a design-time budget: the type of a cipher
  program bounds its orbit depth, and hence its worst-case
  confidentiality loss.

\item \textbf{Expression-tree realization} (\Cref{sec:realizing}).
  We realize the algebra at the program level: a cipher program is
  a composition of cipher maps and plain operations over opaque bit
  strings.
  Cipher-node annotations on a program's expression tree mark which
  subtrees become cipher maps; unmarked nodes run as plain code on
  the cipher bytes.
  The placement of annotations selects the cut points, and the type
  of the cut-point cipher spaces determines the orbit budget per
  \Cref{prop:typed-orbit}.
\end{enumerate}

We deliberately avoid game-based or simulation-based cryptographic
definitions.
Privacy in trapdoor computing comes from the one-way hash and the
representation uniformity of cipher values, not from access-pattern
indistinguishability.
This is not ORAM~\cite{goldreich1996software}, FHE~\cite{gentry2009fully},
or garbled circuits~\cite{yao1982protocols}; it is a distinct paradigm
where totality and noise closure provide the confidentiality guarantees.


%% ====================================================================
\section{Related Work}
\label{sec:related}
%% ====================================================================

\paragraph{Information flow type systems.}
The sum-type impossibility (\Cref{thm:sum-impossibility}) is closely
related to the implicit flow problem in language-based information
security~\cite{sabelfeld2003language}: a conditional branch on a secret
value leaks information through the control flow path taken.
Our result formalizes this for cipher values: pattern matching on a
cipher sum type is an implicit flow that leaks the tag.
The difference is that information flow type systems prevent the leak
at compile time, while cipher types quantify it at the representation
level.

\paragraph{Functional encryption.}
Functional encryption~\cite{boneh2011functional} allows computing
$f(x)$ on an encrypted $x$ by issuing function-specific decryption
keys.
Cipher maps differ in two ways: (1)~they are information-theoretic
(parameterized by $\eta$, $\varepsilon$, $\delta$) rather than
game-based, and (2)~they are approximate (the untrusted machine
evaluates a total function with bounded error, not an exact
decryption).
The type-constructor trade-offs we identify (tag leakage for sums,
correlation leakage for products) apply equally to functional
encryption schemes that support algebraic operations.

\paragraph{Searchable symmetric encryption.}
SSE originates with Song, Wagner, and
Perrig~\cite{song2000practical} and is placed on rigorous
foundations by Curtmola et al.~\cite{curtmola2006searchable}.
The Boolean search construction in \Cref{sec:cipher-bool} provides
an alternative with information-theoretic parameters: cipher sets
with cipher Boolean AND/OR/NOT, rather than encrypted indices with
leakage profiles.
Subsequent work established that leakage profiles in SSE are
attackable in practice: Islam, Kuzu, and
Kantarcioglu~\cite{islam2012access} demonstrated access-pattern
disclosure attacks; Cash et al.~\cite{cash2015leakage} and
Grubbs, Ristenpart, and Shmatikov~\cite{grubbs2017leakage} mounted
leakage-abuse attacks against deployed SSE and encrypted databases;
Naveed, Kamara, and Wright~\cite{naveed2015inference} did the same for
property-preserving encryption.
Cipher maps avoid property preservation entirely and limit the
adversary to the orbit-closure bound of
\Cref{thm:confidentiality-bound}.

\paragraph{Quantitative information flow.}
The entropy-form confidentiality bound of
\Cref{thm:confidentiality-bound} places this work in the quantitative
information flow tradition~\cite{smith2009foundations}, where leakage
is measured as conditional entropy or min-entropy of the secret given
the adversary's observations.
Our contribution is to specialize this framework to cipher values
under an algebraic operation set, yielding a construction-independent
bound in terms of orbit size.

\paragraph{Related concepts.}
Cipher maps realize the oblivious evaluation of a computable function
(after Turing~\cite{turing1936computable}) through a total function on
bit strings, and the cipher Boolean construction
(\Cref{sec:cipher-bool}) is structurally related to the Bloom
filter~\cite{bloom1970space}, which it generalizes by allocating an
explicit False region alongside the True region.
Our experiments sit on the 20~Newsgroups corpus of
Lang~\cite{lang1995newsweeder}.
Maximizing-confidentiality analyses complementary to this paper's
information-theoretic bound appear in a companion
work~\cite{towell2026maxconf}.

%% ====================================================================
\section{Preliminaries}
\label{sec:prelim}
%% ====================================================================

We recall the cipher map abstraction from~\cite{towell2026cipher}.
All notation follows that paper; we cite rather than re-derive.

\begin{definition}[Cipher map~{\cite[Def.~3.1]{towell2026cipher}}]
\label{def:cipher-map}
A \emph{cipher map} for a latent function $f : X \to Y$ is a tuple
$(\fhat, \enc, \dec, s)$ where
$\fhat : \B^n \to \B^n$ is a total function on $n$-bit strings,
$\enc : X \times \{0,\ldots,K{-}1\} \to \B^n$ is an encoding function,
$\dec : \B^n \to Y \cup \{\bot\}$ is a decoding function, and
$s$ is a secret (the trapdoor).
The untrusted machine holds $\fhat$.
The trusted machine holds $\enc$, $\dec$, and $s$.
\end{definition}

A cipher map satisfies four properties, parameterized by
$(\eta, \varepsilon, \mu, \delta)$~\cite[Sec.~4]{towell2026cipher}:

\begin{enumerate}
\item \textbf{Totality.} $\fhat$ is defined on all $2^n$ inputs.
  Out-of-domain outputs are indistinguishable from uniform under the
  random oracle model~\cite{bellare1993random}.

\item \textbf{Representation uniformity} ($\delta$-bounded).
  The marginal distribution of cipher values is $\delta$-close (in total
  variation distance) to uniform over $\B^n$.

\item \textbf{Correctness} ($\eta$-bounded).
  For a random in-domain element $x$ and random representation index $k$,
  $\Pr[\dec(\fhat(\enc(x,k))) \neq f(x)] \leq \eta$.

\item \textbf{Composability.}
  For cipher maps $\fhat$ (correctness $\eta_f$) and $\hat{g}$ (correctness
  $\eta_g$), the composition $\hat{g} \circ \fhat$ has correctness
  $\eta_{g \circ f} \leq 1 - (1-\eta_f)(1-\eta_g)$, with equality under
  the re-randomization condition
  of~\cite[Thm.~4.2]{towell2026cipher}.
\end{enumerate}

The Bernoulli error model~\cite{towell2026bernoulli} provides the
quantitative foundation: element-wise independence of errors reduces the
composition analysis from exponentially many joint failure modes to
per-element Bernoulli parameters.
We use $h : \B^* \to \B^n$ throughout to denote a cryptographic hash
modeled as a random oracle.

A cipher value $\enc(v, k)$ is a cipher map for the constant function
$f_v(x) = v$; we use ``cipher value'' and ``cipher map from a trivial
domain'' interchangeably.

\begin{definition}[Trusted and untrusted machines~{\cite[Sec.~5]{towell2026cipher}}]
\label{def:machines}
The \emph{trusted machine} $T$ holds the trapdoor $s$, $\enc$, $\dec$, and
can encode, decode, and inject filler queries.
The \emph{untrusted machine} $U$ holds $\fhat$ and a collection of cipher
values.
$U$ can evaluate $\fhat$ on any bit string; it cannot decode, distinguish
real from filler queries, or determine the domain $X$.
\end{definition}

\begin{definition}[Cipher type]
\label{def:cipher-type}
For a finite latent type $X$ with encoding $\enc : X \times
\{0, \ldots, K{-}1\} \to \B^n$, the \emph{cipher type} of $X$ is
\[
\cipher{X} \;=\; \{\enc(x, k) : x \in X,\; 0 \leq k < K(x)\} \;\subseteq\; \B^n.
\]
We use $\cipher{X}$ for both the set and the type-theoretic constructor:
$\cipher{X \times Y}$ denotes the cipher type for the latent product, and
similarly for sums and exponentials.
A cipher map $\fhat : \cipher{X} \to \cipher{Y}$ is the cipher form of
the latent function $f : X \to Y$ (extended totally over $\B^n$ per
Property~1).
\end{definition}


%% ====================================================================
\section{Cipher Type Constructors}
\label{sec:types}
%% ====================================================================

Cipher types form an algebra: starting from primitive ciphered values,
the type constructors of ordinary programming (void, unit, product,
sum, exponential) build complex cipher types from simpler ones.
We examine each constructor in turn and identify its
\emph{confidentiality cost}: the irreducible information the untrusted
machine learns when the constructor is realized over cipher values.

A common knob runs through the algebra.
For every binary type constructor, the cipher type can be built by
\emph{joint encoding} (the constructed type as a single cipher value,
$\cipher{A \star B}$) or \emph{component-wise encoding} (the
constructed type as a structure over the cipher types of its
components, $\cipher{A} \star \cipher{B}$).
Joint encoding hides more structure but demands a fresh cipher map for
each operation that exposes it (projection, pattern matching);
component-wise encoding makes operations available to the untrusted
machine at the cost of leaking exactly the structural relationship
that joint encoding was hiding.
The cost is constructor-specific:
products leak \emph{correlations} (\Cref{prop:product-tradeoff}),
sums leak the \emph{tag} (\Cref{thm:sum-impossibility}),
and exponentials leak the \emph{computation pattern}
(\Cref{prop:exponential-tradeoff}).

The same encoding-granularity knob extends to the orbit-closure
analysis of \Cref{sec:orbit}: the operations a constructor exposes
to the untrusted machine determine the size of the reachable cipher
space, and the typed-composition discipline of
\Cref{sec:typed-chains} bounds the resulting confidentiality loss.
Together, the per-constructor costs of this section and the uniform
information-theoretic bound of the next combine into a confidentiality
calculus over the cipher-type algebra.


%% --------------------------------------------------------------------
\subsection{Void and Unit}
\label{sec:void-unit}
%% --------------------------------------------------------------------

The void type $\mathbf{0}$ has no values.
A cipher map $\fhat$ for a function with domain $\mathbf{0}$ is vacuously
correct ($\eta = 0$) and carries no information---it is a total function
on $\B^n$ that is never queried on a valid encoding.
From the untrusted machine's perspective, $\fhat$ is indistinguishable from
a random function.

The unit type $\mathbf{1}$ has exactly one value, $\star$.
A cipher value $\enc(\star, k) \in \B^n$ encodes $\star$; with $K(\star) \geq 1$
representations, all cipher values are representations of the same
underlying value.
The untrusted machine knows that every valid cipher value encodes $\star$
but cannot distinguish valid encodings from noise (by totality).
A cipher map $\fhat : \B^n \to \B^n$ for a function from $\mathbf{1}$ is
a cipher constant: it maps every valid encoding of $\star$ to a cipher
encoding of $f(\star)$.


%% --------------------------------------------------------------------
\subsection{Product Types}
\label{sec:products}
%% --------------------------------------------------------------------

The product type $A \times B$ consists of pairs $(a, b)$ with $a \in A$ and
$b \in B$.
In plaintext computing, a product comes with two projections
$\pi_1 : A \times B \to A$ and $\pi_2 : A \times B \to B$.
Under cipher encoding, we have two choices.

\paragraph{Joint encoding: $\cipher{A \times B}$.}
Encode the pair $(a, b)$ as a single cipher value $\enc((a,b), k) \in \B^n$.
The untrusted machine sees one opaque bit string.
It cannot determine $a$ or $b$ individually, nor can it detect correlations
between $a$ and $b$.
But the projections $\pi_1, \pi_2$ are not available: to extract $a$ from a
cipher encoding of $(a,b)$, one would need a cipher map
$\hat{\pi}_1 : \B^n \to \B^n$ such that
$\dec_A(\hat{\pi}_1(\enc((a,b),k))) = a$---which requires a separate cipher
map construction for each projection.

\paragraph{Component-wise encoding:
$\cipher{A} \times \cipher{B}$.}
Encode $a$ and $b$ independently as
$(\enc_A(a, k_1),\allowbreak \enc_B(b, k_2))$.
The untrusted machine sees a pair of cipher values.
Projections are trivially available (just take the first or second
component).
But the joint distribution of $(a, b)$ is no longer hidden: even if each
component has marginal representation uniformity ($\delta \approx 0$),
the correlation between $a$ and $b$ is preserved in the joint distribution
of the cipher pair.
This is the encoding granularity principle
of~\cite[Sec.~8.1]{towell2026cipher}: coarser encoding hides more
correlations but blocks more operations.

\begin{proposition}[Product confidentiality trade-off]
\label{prop:product-tradeoff}
Let $A$ and $B$ be finite types.
\begin{enumerate}
\item Under joint encoding $\cipher{A \times B}$, the untrusted machine
  cannot distinguish $(a_1, b_1)$ from $(a_2, b_2)$ (up to $\delta$),
  but projections require a cipher map (and thus a separate construction
  by the trusted machine).
\item Under component-wise encoding $\cipher{A} \times \cipher{B}$,
  projections are free, but an adversary observing $N$ pairs
  $(\enc_A(a_i, k_{1i}), \enc_B(b_i, k_{2i}))$ can estimate the joint
  distribution of $(A, B)$ to arbitrary precision as $N \to \infty$.
\end{enumerate}
\end{proposition}

\begin{proof}
Part~(1): Under joint encoding with representation uniformity $\delta$,
each cipher value is $\delta$-close to uniform over $\B^n$.
Without the trapdoor, the untrusted machine cannot decode.
Projections require mapping a cipher encoding of $(a,b)$ to a cipher
encoding of $a$; this is a function from the code space of $A \times B$ to
the code space of $A$, which is exactly a cipher map and requires
construction by the trusted machine.

Part~(2): Under component-wise encoding, the $i$-th observation is
$(\enc_A(a_i, k_{1i}),\allowbreak \enc_B(b_i, k_{2i}))$.
Even with perfect marginal uniformity ($\delta = 0$), each $\enc_A$ is a
deterministic function of $(a_i, k_{1i})$.
If the same representation index is reused, repeated values produce
identical cipher values; even without reuse, the number of distinct cipher
values associated with each latent value is finite ($K(a)$ for each $a$).
Given $N$ observations, the adversary builds the empirical joint distribution
over cipher-value pairs.
Since each cipher value maps to at most one latent value (given a fixed
representation index), the empirical joint distribution of cipher pairs
converges to a mixture of the latent joint distribution.
By standard convergence results, $N = O(|A| \cdot |B| / \xi^2)$
observations suffice to estimate the joint to accuracy $\xi$ in total
variation distance.
\end{proof}


%% --------------------------------------------------------------------
\subsection{Sum Types}
\label{sec:sums}
%% --------------------------------------------------------------------

The sum type $A + B$ is a tagged union: each value is either
$\mathrm{inl}(a)$ for some $a \in A$, or $\mathrm{inr}(b)$ for some
$b \in B$.
The tag (left vs.\ right) enables pattern matching: given a value of type
$A + B$, one inspects the tag to determine which handler to apply.

In plaintext computing, the tag is always available and pattern matching
is free.
Under cipher encoding, this is no longer the case.

\begin{theorem}[Sum-type impossibility]
\label{thm:sum-impossibility}
Let $A$ and $B$ be non-empty finite types with $|A| \geq 1$ and
$|B| \geq 1$.
For cipher values of sum type $A + B$, the following two properties are
mutually exclusive:
\begin{enumerate}
\item \textbf{Tag hiding}: the untrusted machine cannot determine, from a
  cipher value alone, whether it encodes an $A$-value or a $B$-value with
  advantage better than $\delta$ (the representation uniformity parameter).
\item \textbf{Untrusted pattern matching}: holding the $A$-handler
  $\fhat$ and the $B$-handler $\hat{g}$ separately, the untrusted machine
  can apply exactly one of them to a cipher value, choosing which as a
  function of that value, with no per-query assistance from the trusted
  machine.
\end{enumerate}
\end{theorem}

\begin{proof}
We prove the two directions by exhibiting the information-theoretic
constraint.

\medskip
\noindent\textbf{Joint encoding ($\cipher{A+B}$) achieves tag hiding but
blocks pattern matching.}
Under joint encoding, each value $v \in A + B$ (whether $\mathrm{inl}(a)$ or
$\mathrm{inr}(b)$) is encoded as a single cipher value $\enc(v, k) \in \B^n$.
Representation uniformity ($\delta$-bounded) guarantees that the marginal
distribution of $\enc(v, k)$ over random $v$ and $k$ is $\delta$-close to
uniform on $\B^n$.
In particular, the set of cipher values encoding $A$-values and the set
encoding $B$-values are not efficiently distinguishable by the untrusted
machine: both are subsets of $\B^n$, and their union is $\delta$-close to
covering $\B^n$ uniformly.
Therefore the tag (which variant) is hidden up to $\delta$.

Now suppose the untrusted machine could nonetheless perform dispatch:
holding $\fhat$ and $\hat{g}$ separately, it applies exactly one of them
to $c = \enc(v, k)$, choosing which as a function of $c$.
That choice is a selector $s : \B^n \to \{A, B\}$ with
$s(\enc(\mathrm{inl}(a), k)) = A$ for all $a, k$ and
$s(\enc(\mathrm{inr}(b), k)) = B$ for all $b, k$.
But such an $s$ is exactly a distinguisher between $A$-encodings and
$B$-encodings, which exists only if the two sets of cipher values are
separable without the trapdoor.
If $s$ is correct with probability $1 - \gamma$, then the tag is leaked
with advantage at least $1 - \gamma - \delta$, contradicting tag hiding
when $\gamma$ is small.
Note the obligation is to form the \emph{selector}, not merely to compute
the eliminator's result: the result of the case analysis can be produced
as a cipher value with the tag hidden (by a fused eliminator, or by
evaluating both branches), but neither route forms $s$ or constitutes
untrusted dispatch (\Cref{rem:oblivious-elim}).

\medskip
\noindent\textbf{Component-wise encoding ($\cipher{A} + \cipher{B}$)
permits pattern matching but leaks the tag.}
Under component-wise encoding, the value is stored as a tagged pair:
either $(\texttt{L}, \enc_A(a, k))$ or $(\texttt{R}, \enc_B(b, k))$,
where the tag $\texttt{L}$ or $\texttt{R}$ is in the clear (or
distinguishable by the encoding structure---for instance, $A$-cipher-values
and $B$-cipher-values occupy disjoint regions of bit-string space, or
carry a type identifier).
The untrusted machine reads the tag and routes to the correct handler.
Pattern matching is available.

But the tag is observable.
Given a cipher value, the untrusted machine learns whether it encodes an
$A$-value or a $B$-value.
For any distribution $D$ on $A + B$, the fraction of queries that are
$A$-values is observable, leaking $\Pr_{v \sim D}[v \in A]$.

\medskip
\noindent\textbf{No intermediate encoding achieves both.}
Suppose an encoding achieves both tag hiding (advantage $\leq \delta'$
for some small $\delta'$) and untrusted dispatch (success
probability $\geq 1 - \gamma$).
Dispatch requires a selector $s : \B^n \to \{A, B\}$ that is
correct on valid encodings with probability $\geq 1 - \gamma$.
Tag hiding requires that no such function exists with advantage better than
$\delta'$.
Let $p = \Pr[v \in A]$ be the prior probability of an $A$-value.
A trivial strategy that ignores the encoding and always guesses the
majority type achieves accuracy $\max(p, 1-p)$.
The advantage of $s$ over this baseline is at least
$(1 - \gamma) - \max(p, 1-p)$.
Tag hiding requires this advantage to be at most $\delta'$.
Therefore $(1-\gamma) - \max(p, 1-p) \leq \delta'$, which means either
$\gamma \geq 1 - \max(p, 1-p) - \delta'$ (dispatch fails often)
or $\delta' \geq (1-\gamma) - \max(p, 1-p)$ (the tag is substantially
leaked).
For the balanced case $p = 1/2$, this simplifies to
$1/2 - \gamma \leq \delta'$.
Unbalanced distributions ($p \neq 1/2$) make the impossibility stronger:
the baseline accuracy is higher, leaving less room for the encoding to
add information without leaking the tag.
\end{proof}

\begin{remark}[Contrast with products]
\label{rem:sum-vs-product}
Products leak correlations; sums leak the tag.
The asymmetry is structural.
For a product $A \times B$, the two components are always both present;
the question is whether their relationship (correlation) is visible.
For a sum $A + B$, only one component is present; the question is whether
which one is visible.
A product hides \emph{what values co-occur}; a sum hides \emph{what type
is occupied}.
\end{remark}

\begin{remark}[Trusted pattern matching]
\label{rem:trusted-matching}
Under joint encoding $\cipher{A+B}$, pattern matching is still possible---by
the trusted machine.
The trusted machine decodes $\dec(c)$, inspects the tag, applies the
appropriate handler in plaintext, and re-encodes the result.
This requires a round trip through the trusted machine for every case
analysis.
The cost is a communication round, not an impossibility.
The impossibility is for \emph{untrusted} pattern matching without the
trapdoor.
\end{remark}

\begin{remark}[Oblivious elimination without dispatch]
\label{rem:oblivious-elim}
\Cref{thm:sum-impossibility} forbids \emph{dispatch}: selecting which of
two separately held handlers to apply.
It does not forbid computing the \emph{result} of a case analysis as a
cipher value with the tag hidden, which is possible in two ways, neither
of which is untrusted dispatch.
First, the trusted machine can compile the entire eliminator
$h(\mathrm{inl}\,a) = f(a)$, $h(\mathrm{inr}\,b) = g(b)$ into a single
cipher map $\hat{h} : \cipher{A+B} \to \cipher{Z}$ over the joint
encoding; the untrusted machine evaluates $\hat{h}(c)$ blindly, learning
neither the tag nor which branch fired.
This is the efficient form of trusted pattern matching
(\Cref{rem:trusted-matching}): the case analysis is compiled once at
construction time instead of paid as a per-query round trip, and no
trapdoor is needed at evaluation.
Crucially $\hat{h}$ is a trusted construction over the whole domain
$A+B$ (building it requires $f$, $g$, and the secret); the untrusted
machine can evaluate it but could not have formed it, and applying one
fixed map to every input is not the data-dependent selection that
\Cref{thm:sum-impossibility} rules out.
Second, the untrusted machine can evaluate \emph{both} handlers,
$\fhat(c)$ and $\hat{g}(c)$, and obliviously combine them through a
cipher multiplexer (the branching construction of
\Cref{ex:branching-orbit}); this hides the tag as well, at the cost of
evaluating both branches.
In both routes the case analysis is internalized---fused into one map,
or computed on both branches---so the untrusted machine never forms the
selector $s$ and never branches on the tag.
\end{remark}

\begin{example}[Cipher optionals]
\label{ex:optional}
The optional type $\mathrm{Maybe}(A) = \mathbf{1} + A$ is a sum.
Under $\cipher{\mathbf{1} + A}$, the untrusted machine cannot tell
whether a cipher value encodes $\mathrm{Nothing}$ or $\mathrm{Just}(a)$.
Under $\cipher{\mathbf{1}} + \cipher{A}$, the untrusted machine knows
whether the value is present, but not what it is.
For a database application, the first encoding hides which records have null
fields; the second reveals the null pattern but permits null-checking without
the trapdoor.
\end{example}


%% --------------------------------------------------------------------
\subsection{Exponential Types}
\label{sec:exponentials}
%% --------------------------------------------------------------------

The exponential type $B^A = A \to B$ is the type of functions from $A$
to $B$.
The cipher exponential closes the algebra back on itself.

\begin{proposition}[Cipher exponential is the cipher map abstraction]
\label{prop:exponential-identity}
For finite types $A$ and $B$, the cipher exponential $\cipher{A \to B}$
is precisely the set of cipher maps for functions $f : A \to B$.
A function $\fhat : \B^n \to \B^n$ is a well-formed element of
$\cipher{A \to B}$ if and only if it satisfies the four cipher-map
properties of \Cref{def:cipher-map} (totality, representation
uniformity, correctness, composability).
\end{proposition}

\begin{proof}
A cipher map $\fhat$ for $f : A \to B$ is by construction a total
function on $\B^n$ that, when restricted to valid encodings of $A$,
produces approximate encodings of $B$-values
(per~\cite{towell2026cipher}).
Conversely, any total function on $\B^n$ that is correct (with
parameter $\eta$) on some encoding of $A$ is a cipher map for the
function it computes on the latent values.
The four cipher-map properties are exactly the well-formedness
conditions for membership in $\cipher{A \to B}$.
\end{proof}

This identity has a structural consequence: every cipher type
constructor in the paper, including products, sums, and exponentials,
ultimately reduces to a question about cipher maps.
The orbit-closure framework of \Cref{sec:orbit} therefore applies
uniformly to operations on cipher exponentials, just as it applies
to operations on cipher products and sums.

The encoding-granularity knob applies to exponentials too, with its
own constructor-specific cost.

\begin{proposition}[Exponential confidentiality trade-off]
\label{prop:exponential-tradeoff}
Let $f : A \to B$ and $g : C \to D$ be functions with cipher
realizations $\fhat$ and $\hat{g}$.
\begin{enumerate}
\item Under \emph{component-wise} application
  $(\fhat, \hat{g}) : \cipher{A} \times \cipher{C} \to
  \cipher{B} \times \cipher{D}$, the untrusted machine sees the two
  cipher maps applied independently and learns that both $f$ and $g$
  are being computed (the \emph{computation pattern}).
\item Under \emph{joint} application
  $\widehat{(f, g)} : \cipher{A \times C} \to \cipher{B \times D}$
  realized as a single cipher map for the product function
  $(f, g)(a, c) = (f(a), g(c))$, the untrusted machine sees one
  opaque cipher-map invocation and cannot determine which constituent
  functions were applied.
  This requires a dedicated cipher map of size
  $O(|A| \cdot |C|)$ rather than $O(|A|) + O(|C|)$.
\end{enumerate}
\end{proposition}

\begin{proof}
Part~(1): The untrusted machine evaluates $\fhat$ and $\hat{g}$
separately and observes that two distinct cipher maps were applied.
The identities of $\fhat$ and $\hat{g}$, as objects in
$\cipher{A \to B}$ and $\cipher{C \to D}$ respectively, are
themselves part of the cipher-program code visible to the untrusted
machine.

Part~(2): A single cipher map $\widehat{(f, g)}$ over the product
domain has only one invocation, indistinguishable from any other
cipher-map invocation of the same arity and shape.
The product-domain construction has size proportional to
$|A \times C| = |A| \cdot |C|$ rather than the sum of the component
domains, per the lookup-table cost in~\cite[Sec.~6]{towell2026cipher}.
\end{proof}

The asymmetry in part~(2) recapitulates the product-type trade-off
of \Cref{prop:product-tradeoff}: hiding the computation pattern is
itself an instance of joint product encoding, applied to the
exponentials rather than to their inputs.


%% --------------------------------------------------------------------
\subsection{Summary: The Cipher-Type Algebra}
\label{sec:type-summary}
%% --------------------------------------------------------------------

\Cref{tab:algebra-summary} consolidates the per-constructor results.
Three observations stand out.

\begin{table}[ht]
\centering
\small
\caption{Confidentiality cost by type constructor.
``Joint'' encodes the constructed type as a single cipher value;
``component-wise'' builds the cipher type from cipher types of the
components.}
\label{tab:algebra-summary}
\begin{tabular}{@{}lll@{}}
\toprule
Constructor & Joint encoding hides & Component-wise encoding leaks \\
\midrule
$\mathbf{0}$ (void) & N/A & N/A \\
$\mathbf{1}$ (unit) & nothing to hide & nothing to leak \\
$A \times B$ (\Cref{prop:product-tradeoff}) & correlations between $a$, $b$ & joint distribution as $N \to \infty$ \\
$A + B$ (\Cref{thm:sum-impossibility}) & the tag (which variant) & the tag, by impossibility \\
$A \to B$ (\Cref{prop:exponential-tradeoff}) & the computation pattern & which functions are being applied \\
\bottomrule
\end{tabular}
\end{table}

\paragraph{Universal granularity knob.}
For each binary constructor, joint versus component-wise encoding
forms a continuous spectrum (\Cref{rem:sum-vs-product}, the
entanglement parameter $p$ of~\cite[Sec.~8.1]{towell2026cipher}), and
the constructor determines \emph{what} sits on each side.
The product spectrum runs from full correlation hiding to projections
that leak it; the sum spectrum runs from full tag hiding to pattern
matching that leaks it; the exponential spectrum runs from a single
product cipher map to independent component cipher maps.

\paragraph{Sums are categorically different.}
Products and exponentials admit \emph{intermediate} encodings: an
adversary can recover correlations or computation patterns
asymptotically but never exactly from a finite observation.
The sum-type result is sharper: tag hiding and untrusted pattern
matching cannot \emph{coexist} at any encoding, intermediate or
otherwise (\Cref{thm:sum-impossibility}).
The asymmetry is structural, not quantitative: products and
exponentials always have both components present so the question is
which relationships are observable; sums commit to one variant per
value so the question is whether the variant itself is observable.

\paragraph{Closure under the orbit framework.}
The per-constructor costs of \Cref{tab:algebra-summary} are local:
they describe what one cipher value or one operation reveals.
The next section lifts these into a uniform information-theoretic
bound (\Cref{thm:confidentiality-bound}) that controls what the
untrusted machine learns by \emph{actively composing} the operations
exposed to it.
The typed-composition discipline (\Cref{sec:typed-chains}) is the
algebra's design-time budget: the type of a cipher program bounds
how deeply the untrusted machine can chase those compositions.


%% ====================================================================
\section{Orbit Closure and Information Leakage}
\label{sec:orbit}
%% ====================================================================

The type constructors of \Cref{sec:types} describe static confidentiality
properties: what the untrusted machine can learn from the structure of a
single cipher value.
We now develop a tool for analyzing dynamic confidentiality: what the
untrusted machine can learn by actively applying the cipher maps it holds.


%% --------------------------------------------------------------------
\subsection{Definition}
\label{sec:orbit-def}
%% --------------------------------------------------------------------

\begin{definition}[Orbit closure]
\label{def:orbit}
Let $F = \{\fhat_1, \ldots, \fhat_k\}$ be a set of cipher maps
(total functions $\B^n \to \B^n$) available to the untrusted machine,
and let $c \in \B^n$ be a cipher value.
The \emph{orbit closure} of $c$ under $F$ is the smallest set
$\orbitF(c) \subseteq \B^n$ such that:
\begin{enumerate}
\item $c \in \orbitF(c)$, and
\item if $c' \in \orbitF(c)$ and $\fhat_i \in F$, then
  $\fhat_i(c') \in \orbitF(c)$.
\end{enumerate}
Equivalently, $\orbitF(c)$ is the set of all bit strings reachable from
$c$ by composing operations in $F$ in any order and any number of times.
\end{definition}

Since $\B^n$ is finite, $\orbitF(c)$ is always finite and can be computed
by iterative expansion (breadth-first search from $c$, applying each
$\fhat_i$ at each step, until no new values appear).

\begin{remark}
\label{rem:orbit-indep}
The orbit closure depends only on the cipher maps $F$ available to the
untrusted machine and the starting cipher value $c$.
It does not depend on the latent function, the domain, or the trapdoor.
This makes it a tool for analysis by the trusted machine designer: given
the operations you plan to expose, how much can the untrusted machine
explore?
\end{remark}


%% --------------------------------------------------------------------
\subsection{Monotonicity}
\label{sec:monotonicity}
%% --------------------------------------------------------------------

\begin{theorem}[Monotonicity of orbit closure]
\label{thm:monotonicity}
If $F \subseteq G$ (every operation in $F$ is also in $G$), then
$\orbitF(c) \subseteq \mathrm{orbit}_G(c)$ for all $c \in \B^n$.
\end{theorem}

\begin{proof}
By induction on the iterative construction of $\orbitF(c)$.

\emph{Base case.} $c \in \orbitF(c)$ and $c \in \mathrm{orbit}_G(c)$
by definition.

\emph{Inductive step.} Suppose $c' \in \orbitF(c)$ and $c' \in
\mathrm{orbit}_G(c)$ (inductive hypothesis).
For any $\fhat_i \in F$, we have $\fhat_i(c') \in \orbitF(c)$ by
\Cref{def:orbit}.
Since $F \subseteq G$, $\fhat_i \in G$, so
$\fhat_i(c') \in \mathrm{orbit}_G(c)$ by \Cref{def:orbit} applied to $G$.

Every element of $\orbitF(c)$ is therefore also in $\mathrm{orbit}_G(c)$.
\end{proof}

The interpretation is direct: exposing more operations to the untrusted
machine can only enlarge (never shrink) the set of cipher values it can
reach.
Each additional operation is an additional edge in the reachability graph
over $\B^n$.

\begin{corollary}
\label{cor:empty-ops}
$\mathrm{orbit}_\emptyset(c) = \{c\}$ for all $c$.
If the untrusted machine has no cipher maps, it learns nothing beyond the
cipher value it was given.
\end{corollary}

\begin{proof}
With $F = \emptyset$, only the base case of the orbit definition applies.
\end{proof}


%% --------------------------------------------------------------------
\subsection{Confidentiality Bound}
\label{sec:conf-bound}
%% --------------------------------------------------------------------

The orbit closure quantifies reachability.
We now turn it into a confidentiality bound.
The argument is a one-line counting bound on mutual information; the
content is in correctly identifying \emph{what the untrusted machine
observes}.

Let $X$ be the random latent value encoded by $c$, drawn from the
trusted machine's input distribution.
By probing $c$ with the operations in $F$, the untrusted machine
resolves its computation to one of the cipher values reachable from
$c$; write $\mathcal{V}_F(c)$ for this \emph{observation} and model it
as a random variable supported on $\orbitF(c)$.
Two facts pin down its role.
First, the orbit \emph{set} is computable from $c$ and $F$ alone,
without the trapdoor (\Cref{rem:orbit-indep}); it is therefore not
itself a secret-dependent signal, and the orbit is precisely the
\emph{support} of the observation, not a set-valued random variable in
its own right.
Second, what varies with the latent value, and hence carries
information about $X$, is \emph{which} element of that support the
computation resolves to, and there are at most $|\orbitF(c)|$ of those.
The adversary's residual uncertainty about $X$ is the conditional
entropy $H(X \mid \mathcal{V}_F(c))$.

\begin{theorem}[Confidentiality bound, entropy form]
\label{thm:confidentiality-bound}
Let $F$ be a set of cipher maps available to the untrusted machine and
$c$ a cipher value.
Then
\[
H(X \mid \mathcal{V}_F(c)) \;\geq\; H(X) - \log_2 |\orbitF(c)|.
\]
\end{theorem}

\begin{proof}
The observation $\mathcal{V}_F(c)$ is supported on the finite set
$\orbitF(c)$, so it takes at most $|\orbitF(c)|$ distinct values.
The entropy of a random variable supported on a set of size $M$ is at
most $\log_2 M$ (with equality iff it is uniform), hence
$H(\mathcal{V}_F(c)) \leq \log_2 |\orbitF(c)|$.
Mutual information is bounded by the entropy of either argument, so
\[
I(X; \mathcal{V}_F(c)) \;\leq\; H(\mathcal{V}_F(c)) \;\leq\; \log_2 |\orbitF(c)|,
\]
and substituting into
$H(X \mid \mathcal{V}_F(c)) = H(X) - I(X; \mathcal{V}_F(c))$ gives the
claim.
\end{proof}

\begin{remark}[Which scale this bounds]
\label{rem:two-scales}
\Cref{thm:confidentiality-bound} is an \emph{active}-adversary bound:
it limits what the untrusted machine learns by composing the
operations it holds, and the controlling quantity is the orbit size,
not the representation-uniformity parameter $\delta$ (Property~2).
This is a different axis from the \emph{marginal} confidentiality of a
single cipher value, whose leakage is governed by $\delta$ through the
entropy-ratio bound of~\cite{towell2026cipher,towell2026maxconf}.
The two are complementary and do not reduce to each other: a cipher
value can be marginally near-uniform ($\delta \approx 0$) yet leak its
latent value through a large orbit if the exposed operation set is
rich, and conversely a small orbit caps active leakage even when
$\delta$ is not negligible.
No value of $\delta$ bounds the orbit, and no orbit bound improves
$\delta$.
The bound is loose: it ignores the functional relationships between
orbit elements (the adversary knows $c' = \fhat(c)$, which constrains
the possible latent values beyond mere set size) and the algebraic
structure of $F$.
Tighter bounds for specific operations are an open problem.
Equality is approached when orbit elements behave like independent
signals about $X$, e.g., when the cipher maps in $F$ are modeled as
random oracles~\cite{bellare1993random}.
\end{remark}

\begin{remark}[Extremes]
When $|\orbitF(c)| \geq |X|$ the bound is vacuous (the subtracted term
exceeds $H(X)$); in principle a rich enough operation set can identify
every latent value.
When $|\orbitF(c)| = 1$ (no operations available, per
\Cref{cor:empty-ops}), the bound is tight in the strongest sense:
$H(X \mid \mathcal{V}_F(c)) = H(X)$, so the untrusted machine learns
nothing about $X$ beyond what the bare cipher value already revealed.
\end{remark}


%% --------------------------------------------------------------------
\subsection{Examples}
\label{sec:orbit-examples}
%% --------------------------------------------------------------------

\begin{example}[Boolean operations]
\label{ex:boolean-orbit}
Let $X = \{0, 1\}$ (Boolean) with cipher maps $\widehat{\mathrm{AND}}$ and
$\widehat{\mathrm{NOT}}$ available to the untrusted machine.
Given a cipher value $c$ encoding some unknown Boolean $x$, the untrusted
machine can compute:
\begin{align*}
c_1 &= \widehat{\mathrm{NOT}}(c), \\
c_2 &= \widehat{\mathrm{AND}}(c, c_1)
     = \widehat{\mathrm{AND}}(c, \widehat{\mathrm{NOT}}(c)).
\end{align*}

Now, $\mathrm{AND}(x, \mathrm{NOT}(x)) = 0$ for all $x$ (up to
correctness $\eta$).
So $c_2$ is an (approximate) cipher encoding of the constant
$\mathrm{FALSE}$.
If the untrusted machine has a reference cipher value $c_{\mathrm{F}}$
known to encode $\mathrm{FALSE}$ (for instance, from a prior computation
$\widehat{\mathrm{AND}}(c', \widehat{\mathrm{NOT}}(c'))$ for some other
cipher value $c'$), it can compare $c_2$ with $c_{\mathrm{F}}$.

This comparison reveals information: if $c_2 = c_{\mathrm{F}}$
(same representation), the untrusted machine confirms the tautology
and potentially identifies the representation index.
Writing $F_0 = \{\widehat{\mathrm{AND}}, \widehat{\mathrm{NOT}}\}$, the
orbit $\mathrm{orbit}_{F_0}(c)$ includes $c$, $c_1$, $c_2$, and all further
compositions---quickly covering a nontrivial fraction of the cipher space
for small $n$.
\end{example}

\begin{example}[Successor on a bounded type]
\label{ex:successor-orbit}
Let $X = \{0, 1, \ldots, m-1\}$ (integers mod $m$) with a single cipher
map $\widehat{\mathrm{succ}}$ implementing the successor function
$x \mapsto (x+1) \bmod m$.
Given any cipher value $c$ encoding some $x$:
\[
\mathrm{orbit}_{\{\widehat{\mathrm{succ}}\}}(c) =
  \{c,\; \widehat{\mathrm{succ}}(c),\;
   \widehat{\mathrm{succ}}^2(c),\; \ldots,\;
   \widehat{\mathrm{succ}}^{m-1}(c)\}.
\]
After $m$ applications, the orbit returns to $c$ (since successor is
cyclic on $\mathbb{Z}/m\mathbb{Z}$).
The orbit has exactly $m$ elements (assuming $\eta = 0$).
The untrusted machine does not learn which element is
$\enc(0, k)$, but it learns the cycle structure: there exists a cycle of
length $m$ starting from $c$.
This reveals $m = |X|$, the domain cardinality.

By \Cref{thm:confidentiality-bound}, the confidentiality is
$1 - m/2^n$.
For $m \ll 2^n$, this is close to 1: the orbit covers a tiny fraction
of the cipher space.
For $m$ close to $2^n$, confidentiality degrades.
\end{example}

\begin{example}[Branching]
\label{ex:branching-orbit}
Suppose the untrusted machine holds cipher maps $\fhat$ and $\hat{g}$
and evaluates both on a cipher value $c$, obtaining $\fhat(c)$ and
$\hat{g}(c)$.
The orbit of $c$ under $\{\fhat, \hat{g}\}$ includes both results plus all
further compositions.
If $\fhat$ and $\hat{g}$ are the two branches of a conditional
(the ``then'' and ``else'' handlers), both cipher results are visible to
the untrusted machine---even though in plaintext computing, only one
branch would be taken.

This is a form of information leakage specific to trapdoor computing:
the untrusted machine does not know which branch is ``correct,'' but it
sees both outcomes, enlarging the orbit and reducing confidentiality.
Mitigating this requires either (a)~evaluating both branches and having the
trusted machine select the correct result, or (b)~encoding the branch
selection as part of the cipher map (a joint encoding of the conditional).
Strategy (a) doubles the computation; strategy (b) requires a dedicated
cipher map construction.
\end{example}


\begin{remark}[Active probing via Boolean operations]
\label{rem:active-probing}
Exposing cipher Boolean operations creates an active probing risk
beyond passive observation.
Given cipher AND and a cipher value $c$, the adversary computes
$\mathrm{AND}(c, c)$, $\mathrm{AND}(c, c')$ for other values $c'$,
and checks structural consistency.
Over many probes, the adversary partitions cipher values into
equivalence classes corresponding to latent Boolean values.
Each probe enlarges the orbit, and confidentiality degrades per
\Cref{thm:confidentiality-bound}.
This is a concrete instance of orbit closure: Boolean operations are
probing tools that trade functionality for confidentiality.
\end{remark}

\subsection{Typed Composition Chains}
\label{sec:typed-chains}

The active probing risk motivates controlling the orbit by
construction.
The orbit grows because the adversary can feed cipher map outputs back
as inputs, composing operations to reach new cipher values.
A direct mitigation: \emph{prevent self-composition by typing the
cipher spaces}.

\begin{definition}[Typed composition chain]
\label{def:typed-chain}
A \emph{typed composition chain} of depth $k$ is a sequence of
cipher maps $\fhat_0, \ldots, \fhat_{k-1}$ where
$\fhat_i : \cipher{A}_i^{a_i} \to \cipher{A}_{i+1}$ has arity
$a_i \geq 1$ over the cipher space $\cipher{A}_i$, and the cipher
spaces $\cipher{A}_0, \ldots, \cipher{A}_k$ are \emph{distinct}
(disjoint hash value sets, different secrets).
\end{definition}

Since $\cipher{A}_i \neq \cipher{A}_j$ for $i \neq j$, the output of
$\fhat_i$ cannot be fed into $\fhat_j$ for $j \neq i+1$.
The adversary can compose at most $k$ operations in sequence and cannot
loop.

\begin{proposition}[Orbit bound for typed chains]
\label{prop:typed-orbit}
Let $V_0 \subseteq \cipher{A}_0$ be a set of $m$ initial cipher values
available to the adversary, and let $F = \{\fhat_0, \ldots, \fhat_{k-1}\}$
be a typed composition chain of depth $k$ with arities
$a_0, \ldots, a_{k-1}$.
Define $N_0 = m$ and $N_{i+1} = N_i^{a_i}$.
Then
\[
|\orbitF(V_0)| \;\leq\; \sum_{i=0}^{k} N_i,
\]
a finite bound determined entirely by $(k, m, a_0, \ldots, a_{k-1})$
and independent of the cipher-space sizes $|\cipher{A}_i|$.
Two useful special cases:
\begin{itemize}
\item \textbf{Single-value start} ($m = 1$), any arities:
  $|\orbitF(V_0)| \leq 1 + k$.
\item \textbf{Unary chain} ($a_i = 1$), any $m$:
  $|\orbitF(V_0)| \leq m(k+1)$.
\end{itemize}
Substituting $|\orbitF(V_0)| \leq \sum_i N_i$ into
\Cref{thm:confidentiality-bound} bounds the corresponding
confidentiality loss at $\log_2 \sum_i N_i$ bits.
\end{proposition}

\begin{proof}
Partition $\orbitF(V_0)$ by level: $L_0 = V_0 \subseteq \cipher{A}_0$,
and $L_{i+1} \subseteq \cipher{A}_{i+1}$ is the image of $\fhat_i$
applied to all $a_i$-tuples of elements of $L_i$.
Distinct cipher spaces prevent back-references, so every orbit element
lies in exactly one $L_i$.
Since $|L_i| \leq |L_{i-1}|^{a_{i-1}}$, induction gives $|L_i| \leq N_i$,
and therefore $|\orbitF(V_0)| = \sum_i |L_i| \leq \sum_i N_i$.
For $m = 1$, $N_i = 1$ at every level (identical inputs produce one
output for a deterministic cipher map), so the sum is $1+k$.
For $a_i = 1$, $N_i = m$ at every level, so the sum is $m(k+1)$.
\end{proof}

\begin{example}[Depth-limited Boolean search]
\label{ex:typed-bool}
Consider Boolean search with maximum query depth 2 (e.g., $w_1 \land w_2$
but not $w_1 \land w_2 \land w_3$):
\begin{itemize}
\item Cipher sets output to $\cipher{\mathrm{Bool}}_0$.
\item $\mathrm{AND}_0 : \cipher{\mathrm{Bool}}_0 \times
  \cipher{\mathrm{Bool}}_0 \to \cipher{\mathrm{Bool}}_1$
  (binary, $a_0 = 2$).
\item $\mathrm{AND}_1 : \cipher{\mathrm{Bool}}_1 \times
  \cipher{\mathrm{Bool}}_1 \to \cipher{\mathrm{Bool}}_2$
  (binary, $a_1 = 2$).
\item No $\mathrm{AND}_2$ is defined.
\end{itemize}
Starting from a single cipher Boolean $c$ (one membership-query
output), the orbit has size at most $1 + k = 3$: namely
$\{c, \mathrm{AND}_0(c, c), \mathrm{AND}_1(\mathrm{AND}_0(c, c),
\mathrm{AND}_0(c, c))\}$.
With $m$ membership-query outputs as initial values, the bound is
$m + m^2 + m^4$, still finite and controlled by the chain depth.
Contrast with a single shared cipher Boolean space, where the
adversary can self-compose indefinitely and the orbit grows up to
$2^n$.
\end{example}

The general principle: \textbf{the type system controls the orbit}.
By choosing how many cipher space levels to define, the system designer
sets the maximum composition depth and therefore the maximum orbit size
and minimum confidentiality.
This is a compile-time (construction-time) decision with no runtime
cost: the untrusted machine simply does not have cipher maps for
deeper levels.


%% ====================================================================
\section{Realizing Cipher Programs}
\label{sec:realizing}
%% ====================================================================

The type-theoretic results of \Cref{sec:types,sec:orbit} bound what any
cipher-program realization can achieve.
We realize cipher programs as expression trees: a program is a DAG whose
selected nodes are replaced by cipher maps and whose remaining
operations run as plain code on the cipher bytes.
This subsumes iterated constructions (a self-loop is one cut point applied
repeatedly) and reduces to a Turing-machine-style realization in the
degenerate case where the only cipher map is the transition step.
We discuss only the expression-tree case here.


%% --------------------------------------------------------------------
\subsection{Expression-Tree Decomposition}
\label{sec:cipher-programs}
%% --------------------------------------------------------------------

The type-theoretic trade-offs of \Cref{sec:products,sec:sums} have a
direct operational realization: given a program represented as an
expression tree, which subtrees should be replaced by cipher maps?

\begin{definition}[Cipher node annotation]
\label{def:cipher-node}
Given a program $P$ represented as an expression tree, a
\emph{cipher node annotation} marks a subset of nodes $\{n_1, \ldots,
n_k\}$ in the tree.  Each marked node $n_i$, together with its entire
subtree, is replaced by a cipher map: the subtree is evaluated at
construction time (on the trusted machine) over a specified domain, and
the resulting mapping is stored as a cipher map backed by a
perfect hash function~\cite{fredman1984storing}.
Unmarked nodes above the cipher nodes become a \emph{combiner}: itself
a cipher map built over the hash-value outputs of the components.
\end{definition}

The placement of cipher node annotations controls the encoding
granularity:

\begin{itemize}
\item \textbf{Annotation at the root}: the entire program is one cipher
  map.  Maximum confidentiality (no intermediate values exposed), maximum
  space ($O(|X|)$ where $X$ is the input domain of the full program).
\item \textbf{Annotation at the leaves}: each leaf function is a
  separate cipher map.  The combiner chains them.  Intermediate cipher
  values are visible to the untrusted machine, leaking correlations
  between components (as in the product-type trade-off of
  \Cref{sec:products}).
\item \textbf{Annotation at intermediate nodes}: a spectrum between the
  two extremes.  Deeper annotations absorb more computation into cipher
  maps (more confidential, more space); shallower annotations expose
  more structure (less space, more leakage).
\end{itemize}

\begin{remark}[Shared variables and compositional leakage]
\label{rem:shared-vars}
When two cipher nodes $f$ and $g$ share an input variable $x$, the
untrusted machine observes two cipher values $\enc_f(x, k_f)$ and
$\enc_g(x, k_g)$ derived from the same latent $x$.
Even if both encodings are individually $\delta$-close to uniform, the
joint distribution of $(\enc_f(x), \enc_g(x))$ preserves the correlation
(as in \cite[Prop.~8.1]{towell2026cipher}).
The only mitigation is to merge $f$ and $g$ into a single cipher node
(eliminating the shared intermediate) or to inject noise queries that
dilute the correlation signal.
\end{remark}

\paragraph{Implementation via tracing.}
A practical realization uses Python function decorators and runtime
tracing.
Functions marked with a \texttt{@cipher\_node} decorator are
intercepted during a tracing pass over the input domain.
The tracing records which cipher nodes are called, with what arguments,
and what they return.
From this trace, the system:
\begin{enumerate}
\item builds a cipher map for each traced cipher node,
\item builds a combiner cipher map over the hash-output tuples of the
  components,
\item infers the input routing (which program arguments feed which
  cipher node) from the traced call patterns.
\end{enumerate}
Short-circuit evaluation (e.g., Python's \texttt{x if cond else y})
is handled by computing all component cipher maps for every input,
even when the original program skips some branches.
The combiner's domain includes hash tuples for both branches.
The resulting composed cipher program evaluates entirely on the
untrusted machine with no round trips to the trusted machine.


%% --------------------------------------------------------------------
\subsection{Cut-Point Structure}
\label{sec:tm-vs-tree}
%% --------------------------------------------------------------------

\begin{definition}[Cut point]
\label{def:cut-point}
Given a plaintext computation $P$ represented as a directed acyclic
graph with functions at nodes and data flow along edges, a \emph{cut
point} is a node whose subgraph of ancestors is replaced, during
construction, by a single cipher map.
Edges entering the cut point carry plaintext on the trusted machine
and cipher values on the untrusted machine; edges leaving the cut
point carry cipher values throughout.
A \emph{cut-point set} is a collection of cut points such that every
input of $P$ reaches at least one of them; the cipher maps at the cut
points together with the unmarked nodes above them (which evaluate as
plain code on the cipher bytes) compute $P$ on the untrusted machine.
\end{definition}

\paragraph{What leaks above the cut points.}
Exactly what remains plaintext above the cut points is the leakage
profile: the collection of intermediate cipher values flowing between
cipher nodes, plus whatever the unmarked plain operations reveal about
those values (equality checks, branching, dictionary lookups).
Multiple cut points sharing a latent input leak that input's
correlation structure across their outputs, per the product-encoding
trade-off (\Cref{prop:product-tradeoff}).

\paragraph{What bounds the orbit.}
A typed-chain decomposition (distinct cipher spaces between
consecutive cut points) inherits the bound $\sum_i N_i$ of
\Cref{prop:typed-orbit}; for a depth-$k$ tree with binary combinators
starting from $m$ distinct latent inputs, the bound is
$\sum_{i=0}^{k} m^{2^i}$.
Decompositions that reuse a single cipher space across cut points
(equivalently, that admit a self-loop in the cipher-map dependency
graph) fall outside the typed-chain hypothesis and require a separate
argument; for deterministic iterated computations of length $T$, a
single-trajectory bound of $T+1$ orbit elements per starting cipher
value is straightforward but the leakage profile is severe:
every per-step plaintext signal needed to advance the loop (selectors,
indices, halting time) is observable.

\begin{example}[Regex matching as a typed cut-point chain]
\label{ex:regex-cut-points}
Consider matching a fixed regular expression $R$ of size $r$ against a
string $x \in \Sigma^\ell$, with DFA state set $Q_R$ (in the worst
case $|Q_R| = 2^{O(r)}$ via subset construction; for many regexes
much smaller).
The natural typed-chain decomposition places $\ell$ cut points
$\hat{\tau}_0, \ldots, \hat{\tau}_{\ell-1}$ with $\hat{\tau}_i :
\cipher{Q_R}_i \times \cipher{\Sigma} \to \cipher{Q_R}_{i+1}$, each
of space $O(|Q_R| \cdot |\Sigma|)$, total
$O(\ell \cdot |Q_R| \cdot |\Sigma|)$.
The orbit is bounded by $1 + \ell$ from \Cref{prop:typed-orbit}
(single state value at each level), and the leakage above the cut
points is the $(\ell + 1)$-tuple of intermediate cipher state values
$c_0, \ldots, c_\ell$.
The same computation can be expressed as a self-loop on a single
cipher space $\cipher{Q_R \times \Sigma}$ with a single transition
cipher map of space $O(|Q_R|^2 \cdot |\Sigma|)$ (constant in $\ell$),
but the leakage then includes a plaintext (next-state-selector,
direction) pair per step, redacting only the tape contents.
The choice between the typed chain and the self-loop is therefore not
about space (the typed chain is asymptotically smaller per cut point
once $\ell > |Q_R|$) but about \emph{which structure leaks}.
\end{example}


%% ====================================================================
\section{Cipher Boolean Algebra and Evaluation}
\label{sec:cipher-bool}
%% ====================================================================

The cipher Boolean type $\cipher{\mathrm{Bool}}$ specializes the
sum-type trade-off to the most fundamental algebraic structure and
provides the concrete instantiation we evaluate on a real corpus.


%% --------------------------------------------------------------------
\subsection{Cipher Boolean Type}
\label{sec:cipher-bool-type}
%% --------------------------------------------------------------------

\begin{definition}[Cipher Boolean type]
\label{def:cipher-bool}
A cipher Boolean type over $\B^n$ is a partition of the hash space into
three disjoint regions:
\begin{itemize}
\item $T \subset \B^n$: the True region ($|T| / 2^n = p_{\mathrm{true}}$),
\item $F \subset \B^n$: the False region ($|F| / 2^n = p_{\mathrm{false}}$),
\item $N = \B^n \setminus (T \cup F)$: the noise region.
\end{itemize}
The partition is derived from a secret via a pseudorandom permutation.
Without the secret, membership in $T$, $F$, or $N$ is indistinguishable
from random.
\end{definition}

\begin{remark}[Generalization beyond Bool]
\label{rem:general-cipher-type}
\Cref{def:cipher-bool} generalizes to any finite type $Y$.
A cipher type $\cipher{Y}$ partitions $\B^n$ into $|Y|$ value regions
plus a noise region, with $|R_y|/2^n \propto \Pr[Y = y]$
(Shannon-optimal allocation per the source coding
theorem~\cite{shannon1948mathematical}).
Operations $f : Y_1 \times \cdots \times Y_k \to Z$ become cipher maps
over the product of input regions.
The cipher Boolean is the special case $Y = \{\mathrm{True},
\mathrm{False}\}$; the same construction applies to enumerated types,
bounded integers, or any codomain where the output distribution is known.
\end{remark}


%% --------------------------------------------------------------------
\subsection{Noise Unreliability}
\label{sec:noise-unreliability}
%% --------------------------------------------------------------------

The noise region $N$ plays a structural role beyond mere error tolerance.

\begin{proposition}[Noise unreliability]
\label{prop:noise-unreliability}
Let $\mathrm{AND}$, $\mathrm{OR}$, $\mathrm{NOT}$ be cipher maps
implementing Boolean operations over the cipher Boolean type.
If any input is a noise value $n \in N$, the output is \emph{uncorrelated}
with the intended Boolean result: it lands in $T$, $F$, or $N$ according
to the allocation ratios $p_T$, $p_F$, $p_N$, independent of the other
input's value.
\end{proposition}

\begin{proof}
The AND, OR, and NOT cipher maps are constructed over the domain
$T \cup F$.
An input from $N$ is outside this domain.
Since the cipher map is total (backed by a PHF with random noise in
unused slots), a noise input maps to a random slot whose contents decode
to a value in $T$ with probability $p_T$, in $F$ with probability $p_F$,
and in $N$ with probability $p_N$.
For the standard allocation ($p_T = 0.05$, $p_F = 0.90$, $p_N = 0.05$),
a noise input produces False 90\% of the time, True 5\%, and noise 5\%.
Crucially, this output is independent of the other (valid) input:
the cipher map cannot ``correct'' for a noise input because it does not
know which inputs are noise.
\end{proof}

\begin{remark}[Noise prevents output certainty]
The noise region ensures that an adversary observing a sequence of
Boolean cipher map evaluations cannot achieve certainty about any
individual result.
Consider an adversary that repeatedly ANDs a cipher value $c$ with
random cipher values $r_1, r_2, \ldots$.
If $c \in T$ (True), then $\mathrm{AND}(c, r_i)$ produces True when
$r_i \in T$ and False when $r_i \in F$.
But when $r_i \in N$ (5\% probability), the output is noise, which the
adversary cannot distinguish from a legitimate False.
After $k$ ANDs, the adversary has seen some True and some False/noise
results, but cannot determine with certainty whether any given False
result is real or noise-induced.

Dually, an adversary that repeatedly ORs $c$ with random values sees
True for $c \in T$ (always True), but again noise inputs produce
unpredictable outputs.
For $c \in F$, the OR produces True only when $r_i \in T$ (5\%),
False when $r_i \in F$ (90\%), and noise when $r_i \in N$ (5\%).
The noise fraction creates a floor of uncertainty in any statistical
test the adversary performs.
\end{remark}

\paragraph{Practical allocation.}
The cipher Boolean allocation $|T|/2^n = 0.05$, $|F|/2^n = 0.90$,
$|N|/2^n = 0.05$ is motivated by the set membership use case.
A cipher set $\hat{1}_A : \cipher{X} \to \cipher{\mathrm{Bool}}$
maps domain elements to cipher Booleans.
Members are forced into $T$ by the construction.
Non-members land randomly: 90\% decode as False (correct rejection),
5\% as True (false positive), 5\% as noise.
The false positive rate equals $p_{\mathrm{true}}$ and is a system
parameter, not dependent on the set size or universe.

\paragraph{Composition with set operations.}
Boolean search queries like ``$w_1 \land w_2 \land \lnot w_3$'' are
evaluated entirely on the untrusted machine by chaining cipher set
evaluations through cipher AND/OR/NOT maps.
Each Boolean operation is itself a cipher map (constructed via PHF
over the $|T \cup F|^2$ input pairs for AND/OR, or $|T \cup F|$
for NOT).
The untrusted machine sees opaque hash values flowing through opaque
total functions at every stage.


%% --------------------------------------------------------------------
\subsection{Experimental Validation}
\label{sec:evaluation}
%% --------------------------------------------------------------------

A reference implementation (\texttt{cipher-maps}, Python with PHF
backend via \texttt{phobic}) validates the construction on the 20
Newsgroups corpus~\cite{lang1995newsweeder} (18{,}266 documents,
58{,}903 unique words).
With $n = 8$ bits, $|T| = 13$, $|F| = 230$, the AND and OR cipher
maps each have $(13 + 230)^2 = 59{,}049$ entries.
All numbers below are means $\pm$ one standard deviation across five
independent seeds; hardware: x86\_64, single-threaded, Python 3.12.

\paragraph{Boolean search at scale.}
\Cref{tab:boolean-search} shows precision and recall for
representative Boolean queries over a fixed random 5{,}000-document
sample, with all Boolean operations evaluated on the untrusted
machine.
Build rate: 1835 documents per second (mean over seeds).
Plaintext baseline columns are 1.00/1.00 by construction (the cipher
Boolean type is a sound over-approximation of set membership); they
are included to make the cipher-side recall degradation explicit.

\begin{table}[ht]
\centering
\small
\caption{Boolean encrypted search, 5{,}000 documents, $p_T = 0.05$,
$n_\text{seeds} = 5$. Plaintext baseline is 1.00/1.00 by construction.}
\label{tab:boolean-search}
\begin{tabular}{@{}lcccc@{}}
\toprule
Query & Precision & Recall & FP count & Time \\
\midrule
Single term & $0.408 \pm 0.014$ & $1.000 \pm 0.000$ & $233 \pm 13$ & $0.04$s \\
2-term AND & $0.262 \pm 0.024$ & $1.000 \pm 0.000$ & $43 \pm 5$ & $0.13$s \\
3-term AND & $0.291 \pm 0.043$ & $1.000 \pm 0.000$ & $15 \pm 3$ & $0.22$s \\
OR & $0.365 \pm 0.010$ & $0.959 \pm 0.010$ & $437 \pm 16$ & $0.13$s \\
OR AND NOT & $0.330 \pm 0.008$ & $0.909 \pm 0.010$ & $402 \pm 15$ & $0.26$s \\
\bottomrule
\end{tabular}
\end{table}

AND queries maintain perfect recall while false positives decrease
($233 \to 15$ from single term to 3-term AND).
OR and NOT queries lose recall ($0.96$ and $0.91$) due to noise
propagation: when a cipher set returns a noise value $n \in N$,
the subsequent OR or NOT operation produces noise or an unpredictable
result, effectively dropping the document from the result set.
Precision values reflect the $p_T = 0.05$ false-positive rate per
cipher set; for a true set of size $k$ in a corpus of size $N$, the
expected precision is approximately $k / (k + p_T (N - k))$, which
for our representative single-term query yields $\approx 0.4$.

\paragraph{FPR compounding vs.\ Bernoulli model.}
The Bernoulli composition model~\cite{towell2026bernoulli} predicts
that $k$ independent AND operations yield $\mathrm{FPR}_k = p_T^k$
and $k$ independent OR operations yield $1 - (1 - p_T)^k$.
\Cref{tab:fpr-compounding} reports empirical FPRs over $1{,}000$
trials per seed across five seeds.
OR chains match theory tightly across all $k$ (ratio in
$[0.90, 1.02]$).
AND chains diverge sharply with depth: the empirical FPR exceeds
$p_T^k$ by a factor of three at $k = 2$, by $35\times$ at $k = 3$,
and by more than four orders of magnitude at $k = 5$.

\begin{table}[ht]
\centering
\small
\caption{FPR through Boolean chains, 200-document subset,
$1{,}000$ trials per seed, $n_\text{seeds} = 5$.
Last column is empirical / theoretical ratio.}
\label{tab:fpr-compounding}
\begin{tabular}{@{}clrll@{}}
\toprule
Chain & $k$ & Theoretical FPR & Empirical FPR & Ratio \\
\midrule
AND & $1$ & $5.00 \times 10^{-2}$ & $0.0504 \pm 0.0026$ & $1.01$ \\
AND & $2$ & $2.50 \times 10^{-3}$ & $0.0074 \pm 0.0029$ & $2.96$ \\
AND & $3$ & $1.25 \times 10^{-4}$ & $0.0044 \pm 0.0021$ & $35.2$ \\
AND & $4$ & $6.25 \times 10^{-6}$ & $0.0016 \pm 0.0018$ & $256$ \\
AND & $5$ & $3.13 \times 10^{-7}$ & $0.0038 \pm 0.0029$ & $12{,}160$ \\
\midrule
OR & $1$ & $0.0500$ & $0.0448 \pm 0.0086$ & $0.90$ \\
OR & $2$ & $0.0975$ & $0.0912 \pm 0.0078$ & $0.94$ \\
OR & $3$ & $0.1426$ & $0.1452 \pm 0.0051$ & $1.02$ \\
OR & $4$ & $0.1855$ & $0.1724 \pm 0.0101$ & $0.93$ \\
OR & $5$ & $0.2262$ & $0.2098 \pm 0.0062$ & $0.93$ \\
\bottomrule
\end{tabular}
\end{table}

The asymmetry has a structural explanation.
A Bernoulli AND test of $k$ independent draws is itself
$k$-independent, but cipher Boolean AND is a \emph{deterministic
cipher map}: once a chain produces an intermediate cipher Boolean
value $c'$, feeding $c'$ into the next AND produces a fixed (not
redrawn) output.
Independence holds within a single cipher set (each non-member
draws a fresh hash) but breaks when the same cipher Boolean cipher
map is composed with itself $k$ times.
The OR chain sees the same composition pattern but stays close to
theory because the Bernoulli OR formula is dominated by the larger
$1 - (1 - p_T)^k$ term, which is comparatively insensitive to
within-chain dependence; AND collapses to $p_T^k$, a small number
that is easy to overshoot.
A formal correction term is open work; the empirical floor of
$\sim 10^{-3}$ at $k \geq 3$ is one of this paper's concrete
contributions to the cipher-Boolean error model.

\paragraph{Encoding granularity.}
We built a 7-function loan approval pipeline at three granularity
levels (\Cref{tab:granularity}).
The root annotation uses a single PHF-backed cipher map over the
full domain; intermediate and leaf annotations use seed-search-backed
component cipher maps composed via tracing.
Seed-search occasionally fails to find a valid hash assignment for
a given component domain, which is why the intermediate row reports
$n = 2$ successful seeds rather than $5$.

\begin{table}[ht]
\centering
\small
\caption{Encoding granularity over the 150-input loan-approval
domain, $n_\text{seeds} = 5$ attempted (successful seed count
shown for the seed-search levels).}
\label{tab:granularity}
\begin{tabular}{@{}lrrrr@{}}
\toprule
Level & Build (s) & Space & Bits/elem & Intermediates \\
\midrule
Root (1 cipher map) & $0.003 \pm 0.000$ & $694$\,B & $37.0$ & $0$ \\
Intermediate (3 groups, $n = 2$) & $1.67 \pm 1.16$ & $138$\,B & $7.4$ & $3$ \\
Leaf (7 nodes) & $1.14 \pm 1.45$ & $140$\,B & $7.5$ & $7$ \\
\bottomrule
\end{tabular}
\end{table}

As expected from the PHF construction, all three produce zero errors
on the full domain (150 inputs); each cipher map is built over its
declared input set, so on-domain inputs are exact lookups by
construction.
The root annotation is roughly $400\times$ faster (PHF lookup vs.\
per-element seed search) and exposes zero intermediate cipher values.
Counter-intuitively, the leaf and intermediate annotations use
\emph{less} total space ($\approx 140$\,B vs.\ the root's $694$\,B):
each component cipher map's domain is smaller than the product
domain of the full program, and the seed-search backend stores one
$n$-bit seed per element.
The space-leakage trade-off is therefore not a uniform ``coarser is
larger''; coarser annotation buys hidden intermediate values at the
cost of construction time and a build that requires the full input
domain up front, while leaf annotations build smaller cipher maps
that compose dynamically but expose every intermediate hash flow.


%% ====================================================================
\section{Discussion}
\label{sec:discussion}
%% ====================================================================

\paragraph{Relationship to cipher maps.}
This paper extends the cipher map framework of~\cite{towell2026cipher} in
two directions: inward (what happens when cipher values carry algebraic
structure) and outward (what happens when the untrusted machine actively
explores the cipher space).
The cipher map paper focuses on individual maps and their four properties;
this paper studies how those properties interact with type constructors and
how confidentiality degrades under active use.
The two frameworks measure confidentiality on complementary axes.
The cipher map paper bounds the leakage of a single cipher value's
\emph{marginal} distribution through the representation-uniformity
parameter $\delta$~\cite[Prop.~5.1]{towell2026cipher}, and analyzes
cross-instance \emph{collision} leakage when many cipher values are
observed together~\cite[\S 8.2]{towell2026cipher}; the entropy-ratio
treatment is developed further in~\cite{towell2026maxconf}.
Our orbit-closure bound (\Cref{thm:confidentiality-bound}) measures a
third, orthogonal axis: the \emph{reachability} leakage an adversary
accrues by actively composing the operations it holds.
The three are complementary, not competing: a cipher value can be
marginally uniform ($\delta \approx 0$), collision-free, and still
leak through a large orbit if the exposed operation set is rich.

\paragraph{Relationship to Bernoulli data types.}
The accuracy-side analysis of algebraic types over approximate values is
developed in~\cite{towell2026bernoulli}, which studies the Kronecker
factorization of joint confusion matrices for product types and the error
propagation through sum-type constructors.
That paper asks: how does $\eta$ change when types are composed?
This paper asks the complementary question: how does confidentiality
change when types are composed?
Together, they characterize the full trade-off space (accuracy $\times$
confidentiality) for algebraic types over trapdoor computing.

\paragraph{The encoding granularity principle.}
The product and sum trade-offs in \Cref{sec:products,sec:sums} are
instances of a single principle developed
in~\cite[Sec.~8.1]{towell2026cipher}: coarser encoding granularity
(encoding more values as a single unit) yields better confidentiality but
fewer operations available to the untrusted machine.
The entanglement parameter $p$~\cite[Sec.~8.1]{towell2026cipher} controls
the spectrum: $p = 1$ (component-wise) gives maximum functionality and
minimum confidentiality; $p = k$ (whole-state) gives maximum
confidentiality and minimum functionality.

\paragraph{Open questions.}

\begin{enumerate}
\item \emph{Recursive types.}
  Lists, trees, and other recursive types are built from sums and products.
  The sum-type impossibility (\Cref{thm:sum-impossibility}) implies that
  recursive traversal (which requires pattern matching at each node)
  either leaks the structure (which constructors are used at each level)
  or requires a trusted-machine round trip at each node.
  Quantifying the total leakage for a recursive traversal of depth $d$
  is open.

\item \emph{Tight orbit bounds for specific constructions.}
  \Cref{thm:confidentiality-bound} gives a general upper bound on
  confidentiality in terms of orbit size.
  For specific cipher map constructions (e.g., the trapdoor boolean
  algebra of~\cite{towell2026cipher}), tighter bounds may be possible
  by exploiting the structure of the construction.

\item \emph{Plain operations on cipher bytes.}
  The expression-tree realization conflates two roles for the
  unmarked nodes above the cut points: a cipher-map combiner (built
  at construction time over the product of cipher hash spaces) and
  plain code that runs on the cipher bytes at evaluation time.
  Characterizing what each plain operation leaks (equality checks
  reveal collision structure; dictionary lookups reveal hash bucket
  membership) and when the trade with construction-time cost is
  preferable to a precomputed combiner cipher map is open.

\item \emph{Partial sum-type hiding.}
  \Cref{thm:sum-impossibility} shows that full tag hiding and full
  untrusted pattern matching are mutually exclusive.
  Is there a meaningful intermediate notion---hiding the tag from
  some adversary classes while permitting restricted pattern matching?
  For instance, one could imagine a construction where the tag is hidden
  against passive adversaries (who only observe cipher values) but
  available to active adversaries (who evaluate cipher maps).

\item \emph{Cut-point placement as an optimization.}
  \Cref{sec:tm-vs-tree} describes cut points as a design knob over the
  computation DAG.
  Formulating cut-point selection as a constrained optimization,
  subject to the orbit bound of \Cref{prop:typed-orbit} and
  construction-time space budgets, would turn the choice into a
  solvable problem rather than a designer's judgement call.
\end{enumerate}


%% ====================================================================
\section{Conclusion}
%% ====================================================================

Cipher types form an algebra with a coherent confidentiality
calculus.
We surveyed the algebra constructor by constructor: products leak
correlations, sums force a categorical impossibility (tag hiding
and untrusted pattern matching cannot coexist), and the cipher
exponential is the cipher map abstraction itself, closing the
algebra back on the trapdoor framework.
Across the algebra, a uniform information-theoretic bound holds:
the conditional entropy of the latent value given the untrusted
machine's view is at least $H(X) - \log_2 |\mathrm{orbit}|$.
A typed-composition discipline turns this bound into a design-time
budget; these results apply to any cipher map construction, not
just specific implementations.

The algebra realizes at the program level as expression-tree
decomposition.
A cipher program is a composition of cipher maps and plain code over
opaque bit strings: cipher-node annotations on a Python program
select the cut points where cipher maps replace plaintext subtrees,
while unmarked nodes run as plain code on the cipher bytes.
The placement of annotations selects which intermediate structure
the untrusted machine observes; the type discipline of the
cut-point cipher spaces bounds the orbit per
\Cref{prop:typed-orbit}.

On the practical side, cipher Boolean types with AND/OR/NOT as cipher
maps enable Boolean search queries evaluated entirely on the untrusted
machine.
Experiments on the 20 Newsgroups corpus (5{,}000 documents, five
seeds) show perfect recall for AND queries and noise-driven recall
loss to $\approx 0.91$ for queries containing OR or NOT.
The most striking empirical finding is that AND-chain FPR diverges
from the Bernoulli independence prediction by orders of magnitude:
empirical FPR exceeds $p_T^k$ by $35\times$ at $k = 3$ and by
over $10^4\times$ at $k = 5$, while OR chains stay within $10\%$ of
theory at every depth.
This points to a structural correction term for composed deterministic
cipher maps as future work.
The encoding granularity spectrum is measurable: root annotation
(one cipher map for the entire program) is roughly $400\times$ faster
than leaf annotation (one cipher map per function), exposes no
intermediates, but consumes more space than the smaller component
maps; the trade-off is not a single axis but a multi-dimensional
choice over (build time, total space, intermediate exposure).

%% ====================================================================
\section*{Acknowledgments}
%% ====================================================================

The algebraic cipher types concept originates in a 2019--2022 C++
notebook by the author, which explored cipher Booleans, cipher
unions, and several iterated constructions through concrete
implementations.
The present paper extracts and formalizes the theoretical content
from that notebook, building on the cipher map framework developed
in~\cite{towell2026cipher}.

\bibliographystyle{plainnat}
\bibliography{references}

\end{document}
